\chapter{VII7 nel modo maggiore e II7 nel modo minore: stessi suoni, funzioni diverse}
\label{chap:vii7_ii7}

Uno degli aspetti più affascinanti dell'armonia tonale è la capacità di uno stesso insieme di note di assumere significati e funzioni completamente diverse a seconda del contesto tonale. In questo capitolo esploriamo due accordi di settima che, pur essendo costituiti dagli stessi suoni (in tonalità relative), svolgono ruoli armonici distinti e complementari: la \textbf{settima di sensibile} (VII7) nel modo maggiore e la \textbf{settima sul secondo grado} (II7) nel modo minore.

\section{La settima di sensibile nel modo maggiore (VII7)}
\label{sec:vii7_maggiore}

Nel modo maggiore, costruendo un accordo di settima sul settimo grado della scala, otteniamo una quadriade di grande interesse armonico.

\begin{definizione}[Settima di sensibile (VII7)]
	\label{def:vii7}
	La \textbf{settima di sensibile} è un accordo di settima costruito sul VII grado della scala maggiore. È composto da una triade diminuita (fondamentale, terza minore, quinta diminuita) con l'aggiunta di una settima minore. In Do maggiore: Si-Re-Fa-La.
	
	Questo accordo è anche chiamato \textbf{semidiminuito} o \textbf{settima di seconda specie}, indicato con il simbolo VII\textsuperscript{\o 7} o VII7 con la cifratura {\o}.
\end{definizione}

La caratteristica fondamentale di questo accordo è che contiene tre note della settima di dominante, del quale manca la fondamentale. È naturale dunque che questo accordo abbia funzione di dominante e risolva sulla tonica.

\begin{regola}[Funzione e risoluzione della VII7]
	\label{reg:funzione_vii7}
	La VII7 ha funzione di \textbf{dominante}. La sua risoluzione naturale è sulla tonica (I). Secondo le regole generali di risoluzione della dissonanza (Capitolo \ref{chap:teoria_dissonanza}), la settima deve scendere di grado. Le due note che formano il tritono (Si-Fa in Do maggiore) seguono il loro movimento obbligato:
	\begin{itemize}
		\item La sensibile (Si) sale alla tonica (Do)
		\item La quinta diminuita (Fa) scende alla terza della tonica (Mi)
		\item La settima (La) scende anch'essa di grado (Sol)
	\end{itemize}
\end{regola}

\begin{esempio}[VII7 in Do Maggiore e sua risoluzione]
	\label{ex:vii7_maggiore}
	L'accordo di settima di sensibile in stato fondamentale risolve sulla tonica.
	\musica{
		\partilate
		{ \key do \major \time 4/4 }
		{ la'2 sol'2 }
		{ fa'2 mi'2 }
		{ re'2 do'2 }
		{ si,2 do2 }
		\lyricmode { "VIIø7"2 "I"2 }
		\figuremode { <7>2 s2 }
	}
\end{esempio}

\subsection{I rivolti della VII7 e la Regola dell'Ottava}
\label{subsec:rivolti_vii7}

Come tutti gli accordi di settima, anche la VII7 può essere rivoltata. Il suo uso più comune è nel \textbf{primo rivolto} (VII6/5), che appare nella Regola dell'Ottava.

\begin{regola}[VII6/5 nella Regola dell'Ottava discendente]
	\label{reg:vii65_regola}
	Nella scala discendente del modo maggiore, quando il basso raggiunge il VII grado, si utilizza il primo rivolto della settima di sensibile (VII6/5). Questo crea una linea del basso fluida e permette alla settima (che si trova ora al soprano o in una voce intermedia) di risolvere regolarmente scendendo di grado.
\end{regola}

\begin{esempio}[VII6/5 nella Regola dell'Ottava - Do Maggiore]
	\label{ex:vii65_regola_maggiore}
	Frammento della scala discendente con evidenziazione del VII6/5.
	\musica{
		\partilate
		{ \key do \major \time 4/4 }
		{ do''4 si'4 la'4 sol'4 }
		{ mi'4 re'4 mi'4 re'4 }
		{ sol4 sol4 do'4 si4 }
		{ do4 si,4 la,4 sol,4 }
		\lyricmode { "I"4 "VII"4 "VI"4 "V"4 }
		\figuremode { \unset FiguredBass.useBassFigureExtenders
			s4 <6 5>4 <6>4 s4 }
	}
\end{esempio}

\section{La settima sul secondo grado nel modo minore (II7)}
\label{sec:ii7_minore}

Nel modo minore armonico, l'accordo di settima costruito sul secondo grado della scala ha la stessa struttura della VII7 del modo maggiore: una triade diminuita con settima minore (accordo semidiminuito).

\begin{definizione}[Settima sul secondo grado (II7)]
	\label{def:ii7_minore}
	La \textbf{settima sul secondo grado} nel modo minore armonico è un accordo semidiminuito che svolge funzione di \textbf{sottodominante}. In La minore: Si-Re-Fa-La (gli stessi suoni della VII7 di Do maggiore!).
	
	Si indica come II\textsuperscript{\o 7} o semplicemente II7.
\end{definizione}

Tale accordo, come visto nel paragrafo \ref{def:vii7}, contiene tre di quattro note della settima di dominante del modo relativo maggiore e porta alla sua tonica; nel contesto del modo minore, la tonica del modo maggiore corrisponde alla dominante del modo minore: l'accordo II7 nel modo minore, dunque, risolve sulla dominante.

\subsection{La progressione I - II7 - V - I nel modo minore}
\label{subsec:progressione_ii7}

L'uso più efficace della II7 nel modo minore è all'interno della progressione armonica fondamentale che esprime il ciclo Tonica-Sottodominante-Dominante-Tonica.

\begin{regola}[Preparazione della settima sul II grado]
	\label{reg:preparazione_ii7}
	La settima dell'accordo II7, secondo i principi generali di preparazione della dissonanza (cfr. Capitolo \ref{chap:teoria_dissonanza}), deve essere preparata come consonanza nell'accordo precedente (I). La condotta delle voci più fluida si ottiene utilizzando il \textbf{primo rivolto} dell'accordo II7, ponendo la terza al basso (Re).
\end{regola}

\begin{esempio}[Progressione I - II6/5 - V - I in La minore]
	\label{ex:progressione_ii7_minore}
	La settima (La) viene preparata nell'accordo di tonica e risolve scendendo di grado (Sol) nell'accordo di dominante.
	\musica{
		\partilate
		{ \key la \minor \time 4/4 }
		{ do'2 si2 | si1 }
		{ la2 la2 | sold1 }
		{ mi2 fa2 | mi1 }
		{ la,2 re2 | mi2 mi,2 | la,1 }
		\lyricmode { "I"2 "IIø7"2 | "V"1 }
		\figuremode { s2 <6 5>2 | s1 }
	}
\end{esempio}

\subsection{II6/5 nella Regola dell'Ottava (modo minore)}
\label{subsec:ii7_regola_minore}

Analogamente al modo maggiore, anche nel modo minore il secondo grado della scala ascendente può essere armonizzato con il primo rivolto della settima (II6/5).

\begin{regola}[II6/5 nella Regola dell'Ottava ascendente - Modo Minore]
	\label{reg:ii65_regola_minore}
	Nella scala ascendente del modo minore, il II grado al basso viene armonizzato con l'accordo II6/5. La settima deve essere preparata e risolve regolarmente scendendo di grado nell'accordo successivo (III o V).
\end{regola}

\begin{esempio}[II6/5 nella Regola dell'Ottava - La minore]
	\label{ex:ii65_regola_minore}
	Frammento della scala ascendente in La minore.
	\musica{
		\partilate
		{ \key la \minor \time 4/4 }
		{ do'4 re'4 mi'4 fa'4 }
		{ la4 la4 si4 la4 }
		{ mi4 fa4 sold4 do'4 }
		{ la,4 si,4 do4 re4 }
		\lyricmode { "I"4 "II"4 "III"4 "IV"4 }
		\figuremode { s4 <6 5>4 <6>4 s4 }
	}
\end{esempio}

\section{Sintesi: VII7 (Magg) vs II7 (Min)}
\label{sec:sintesi_vii_ii}

\begin{center}
	\begin{tabular}{|l|c|c|}
		\hline
		\textbf{Proprietà} & \textbf{VII7 (Modo Maggiore)} & \textbf{II7 (Modo Minore)} \\
		\hline
		Struttura & Semidiminuita (m-m-M) & Semidiminuita (m-m-M) \\
		Funzione & Dominante & Sottodominante \\
		Risoluzione naturale & Tonica (I) & Dominante (V) \\
		Preparazione settima & Preferibile & Obbligatoria \\
		Note (es. Do M / La m) & Si-Re-Fa-La & Si-Re-Fa-La \\
		\hline
	\end{tabular}
\end{center}
